\appendix

\chapter{Complete Monument Catalog}
\addcontentsline{toc}{chapter}{\texorpdfstring{\MakeUppercase{Appendix A: Monument Catalog}}{Appendix A: Monument Catalog}}

The full catalog of twenty procedural monuments shipped with the system is reproduced below. Every entry is a row in \texttt{WorldMonumentLibrary.\_catalog} and is rendered identically by the Unity and Three.js builds.

\begin{table}[h!]
\centering
\renewcommand{\arraystretch}{1.3}
\setlength{\tabcolsep}{4pt}
\footnotesize
\begin{tabular}{|p{3.4cm}|p{3.6cm}|p{3.0cm}|p{3.4cm}|}
\hline
\textbf{Identifier} & \textbf{Title} & \textbf{Era} & \textbf{Location} \\
\hline
TajMahal           & Taj Mahal                      & 17th C. (Mughal)         & Agra, India \\\hline
QutubMinar         & Qutub Minar                    & 12th C. (Sultanate)      & Delhi, India \\\hline
HampiChariot       & Stone Chariot, Hampi           & 16th C. (Vijayanagara)   & Hampi, India \\\hline
KonarkSunTemple    & Konark Sun Temple              & 13th C. (Eastern Ganga)  & Odisha, India \\\hline
AjantaCaves        & Ajanta Caves                   & 2nd C. BCE -- 5th C. CE  & Maharashtra, India \\\hline
GreatPyramidGiza   & Great Pyramid of Giza          & c. 2560 BCE              & Giza, Egypt \\\hline
Colosseum          & Roman Colosseum                & 80 CE                    & Rome, Italy \\\hline
Stonehenge         & Stonehenge                     & c. 3000 BCE              & Wiltshire, England \\\hline
Parthenon          & Parthenon                      & 447 BCE                  & Athens, Greece \\\hline
EiffelTower        & Eiffel Tower                   & 1889 CE                  & Paris, France \\\hline
ChristTheRedeemer  & Christ the Redeemer            & 1931 CE                  & Rio de Janeiro, Brazil \\\hline
StatueOfLiberty    & Statue of Liberty              & 1886 CE                  & New York, USA \\\hline
EasterIslandMoai   & Moai (Easter Island)           & 1250--1500 CE            & Rapa Nui, Chile \\\hline
AngkorWat          & Angkor Wat                     & 12th C. (Khmer)          & Siem Reap, Cambodia \\\hline
GreatWallSection   & Great Wall (section)           & 7th C. BCE -- 17th C. CE & Northern China \\\hline
PetraTreasury      & Al-Khazneh, Petra              & 1st C. CE (Nabataean)    & Petra, Jordan \\\hline
ChichenItza        & El Castillo, Chich\'en Itz\'a  & 9th--12th C. (Maya)      & Yucat\'an, Mexico \\\hline
LeaningTowerPisa   & Leaning Tower of Pisa          & 1372 CE                  & Pisa, Italy \\\hline
SydneyOperaHouse   & Sydney Opera House             & 1973 CE                  & Sydney, Australia \\\hline
\end{tabular}
\caption{Complete monument catalog with identifier, title, era and location.}
\label{tab:catalog}
\end{table}

\begin{figure}[h]
\centering
\fbox{\rule{0pt}{8cm}\rule{0.85\textwidth}{0pt}}
\caption{Monument silhouette grid: all twenty procedural monuments rendered side by side. Insert montage here.}
\label{fig:montage}
\end{figure}

\chapter{Path Graph Edge Listing}
\addcontentsline{toc}{chapter}{\texorpdfstring{\MakeUppercase{Appendix B: Path Graph}}{Appendix B: Path Graph}}

The full set of path edges installed by \texttt{InstallDefaultMuseumGraph()} is listed below. Each edge has a unique identifier, a source node, a target node, an associated theme and a locomotion mode. Edges marked Locked require the visitor to first collect the Treasury Key from the Archaeological Shards room.

\begin{table}[h!]
\centering
\renewcommand{\arraystretch}{1.3}
\setlength{\tabcolsep}{4pt}
\footnotesize
\begin{tabular}{|p{4.0cm}|p{3.0cm}|p{3.0cm}|p{2.0cm}|p{1.5cm}|}
\hline
\textbf{Edge ID} & \textbf{From} & \textbf{To} & \textbf{Loco.} & \textbf{Locked} \\
\hline
e\_entrance\_main      & entrance         & main\_corridor     & Walking  & No \\\hline
e\_main\_iraq          & main\_corridor   & iraq\_room         & Teleport & No \\\hline
e\_main\_prague        & main\_corridor   & prague\_room       & Teleport & No \\\hline
e\_main\_shards        & main\_corridor   & shards\_room       & Teleport & No \\\hline
e\_main\_aerial        & main\_corridor   & aerial\_room       & Teleport & No \\\hline
e\_main\_indian        & main\_corridor   & indian\_hall       & Teleport & No \\\hline
e\_main\_east          & main\_corridor   & east\_corridor     & Walking  & No \\\hline
e\_east\_colosseum     & east\_corridor   & colosseum\_room    & Teleport & No \\\hline
e\_east\_greece        & east\_corridor   & greece\_room       & Teleport & No \\\hline
e\_colosseum\_aqueduct & colosseum\_room  & roman\_aqueduct    & Walking  & No \\\hline
e\_greece\_inner       & greece\_room     & parthenon\_inner   & Walking  & No \\\hline
e\_aerial\_sky         & aerial\_room     & sky\_observatory   & Floating & No \\\hline
e\_indian\_taj         & indian\_hall     & taj\_mahal\_garden & Walking  & No \\\hline
e\_indian\_konark      & indian\_hall     & konark\_courtyard  & Walking  & No \\\hline
e\_indian\_treasury    & indian\_hall     & hidden\_treasury   & Glide    & Yes \\\hline
e\_shortcut\_iraq\_shards & iraq\_room    & shards\_room       & Walking  & No \\\hline
e\_taj\_konark\_loop   & taj\_mahal\_garden & konark\_courtyard & Walking & No \\\hline
e\_konark\_treasury    & konark\_courtyard & hidden\_treasury  & Walking  & Yes \\\hline
e\_iraq\_aerial        & iraq\_room       & aerial\_room       & Walking  & No \\\hline
\end{tabular}
\caption{Complete path graph edge listing.}
\label{tab:edges}
\end{table}

\chapter{Selected Source Listings}
\addcontentsline{toc}{chapter}{\texorpdfstring{\MakeUppercase{Appendix C: Source Listings}}{Appendix C: Source Listings}}

The following abbreviated source listings show the structural backbone of the system. The complete project comprises approximately ten thousand lines of code; only representative snippets are reproduced here.

\section*{C.1 Monument Builder Dispatch}

\begin{lstlisting}[language={[Sharp]C}]
public static GameObject Build(WorldMonument m, Material mat, Transform parent) {
    var def = GetDefinition(m);
    var root = new GameObject($"Monument_{m}");
    if (parent != null) root.transform.SetParent(parent, false);
    var material = mat ?? CreateDefaultMaterial(def.stoneColor);

    switch (m) {
        case WorldMonument.TajMahal:          BuildTajMahal(root, material); break;
        case WorldMonument.Colosseum:         BuildColosseum(root, material); break;
        case WorldMonument.Stonehenge:        BuildStonehenge(root, material); break;
        case WorldMonument.Parthenon:         BuildParthenon(root, material); break;
        case WorldMonument.EiffelTower:       BuildEiffelTower(root, material); break;
        case WorldMonument.ChristTheRedeemer: BuildChristTheRedeemer(root, material); break;
        case WorldMonument.StatueOfLiberty:   BuildStatueOfLiberty(root, material); break;
        case WorldMonument.EasterIslandMoai:  BuildMoai(root, material); break;
        case WorldMonument.PetraTreasury:     BuildPetraTreasury(root, material); break;
        // ... twelve more monuments ...
    }
    foreach (var r in root.GetComponentsInChildren<MeshRenderer>())
        if (r.sharedMaterial == null) r.sharedMaterial = material;
    root.transform.localScale = Vector3.one * 0.05f * def.scaleMultiplier;
    return root;
}
\end{lstlisting}

\section*{C.2 Path Edge Definition}

\begin{lstlisting}[language={[Sharp]C}]
[Serializable]
public class PathEdge {
    public string edgeId;
    public string fromNodeId;
    public string toNodeId;
    public bool bidirectional = true;
    public PathTheme theme = PathTheme.Default;
    public LocomotionMode locomotion = LocomotionMode.Teleport;
    [TextArea(1, 3)] public string narratorCue;
    public bool requiresUnlock;
    public string unlockEventId;
    public float lengthMeters = 4f;
}
\end{lstlisting}

\section*{C.3 Three.js Animation Loop}

\begin{lstlisting}[language=JavaScript]
renderer.setAnimationLoop(() => {
    const dt = Math.min(clock.getDelta(), 0.05);
    update(dt);              // movement, collisions, room detection
    renderer.render(scene, camera);
});

function update(dt) {
    const speed = (keys.ShiftLeft || keys.ShiftRight) ? player.runSpeed : player.speed;
    const fwd   = new THREE.Vector3(-Math.sin(player.yaw), 0, -Math.cos(player.yaw));
    const right = new THREE.Vector3( Math.cos(player.yaw), 0, -Math.sin(player.yaw));
    const move  = new THREE.Vector3();
    if (keys.KeyW) move.add(fwd);
    if (keys.KeyS) move.sub(fwd);
    if (keys.KeyD) move.add(right);
    if (keys.KeyA) move.sub(right);
    if (move.lengthSq() > 0) move.normalize().multiplyScalar(speed * dt);

    const next = player.pos.clone().add(new THREE.Vector3(move.x, 0, 0));
    if (!collides(next)) player.pos.x = next.x;
    const next2 = player.pos.clone().add(new THREE.Vector3(0, 0, move.z));
    if (!collides(next2)) player.pos.z = next2.z;

    camera.position.copy(player.pos);
    camera.rotation.set(player.pitch, player.yaw, 0, 'YXZ');
    detectRoom();
}
\end{lstlisting}

\section*{C.4 Default Path Graph Installation}

\begin{lstlisting}[language={[Sharp]C}]
public void InstallDefaultMuseumGraph() {
    Clear();
    AddN("entrance",        "Grand Entrance",       MuseumRoom.Corridor);
    AddN("main_corridor",   "Main Corridor",        MuseumRoom.Corridor);
    AddN("iraq_room",       "Iraq -- Nahum Shrine", MuseumRoom.IraqSection);
    AddN("indian_hall",     "Indian Hall",          MuseumRoom.IndianArchitecture);
    AddN("hidden_treasury", "Hidden Treasury",      MuseumRoom.IndianArchitecture);
    // ... remaining nodes ...

    AddE("e_entrance_main",   "entrance",      "main_corridor",   PathTheme.Reflective, LocomotionMode.Walking, 6f);
    AddE("e_main_indian",     "main_corridor", "indian_hall",     PathTheme.Mystical,   LocomotionMode.Teleport, 5f);
    AddE("e_indian_treasury", "indian_hall",   "hidden_treasury", PathTheme.Mystical,   LocomotionMode.Glide, 4f, true, "found_key");
    // ... remaining edges ...
}
\end{lstlisting}

\chapter{Build and Run Instructions}
\addcontentsline{toc}{chapter}{\texorpdfstring{\MakeUppercase{Appendix D: Build and Run}}{Appendix D: Build and Run}}

\section*{D.1 Web Build}

The browser build is the simplest path to experiencing the museum.

\begin{enumerate}
\item Locate \texttt{VirtualMuseum.html} in the project root.
\item Double-click the file. Any modern browser will open it directly.
\item Click \textbf{Enter Museum} on the title screen.
\item If a Meta Quest browser detects WebXR support, an \textbf{Enter VR} button appears at the top of the screen.
\end{enumerate}

\section*{D.2 Unity Build for Meta Quest}

\begin{enumerate}
\item Install Unity Hub and add Unity 2022.3 LTS with the Android Build Support module.
\item Open the project at the repository root with Unity Hub.
\item Allow Unity to resolve packages from \texttt{Packages/manifest.json}.
\item Run \textbf{Tools $\rightarrow$ Virtual Museum $\rightarrow$ Setup Wizard} to validate the scene.
\item In \textbf{File $\rightarrow$ Build Settings}, switch the platform to Android.
\item Connect a Meta Quest in developer mode over USB-C.
\item Click \textbf{Build And Run}.
\end{enumerate}

\section*{D.3 Editor Tools}

The project ships with two custom editor windows:

\begin{itemize}
\item \textbf{Tools $\rightarrow$ Virtual Museum $\rightarrow$ Build Monument Catalog} --- creates ScriptableObject assets and prefabs for the Indian heritage subset.
\item \textbf{Tools $\rightarrow$ Virtual Museum $\rightarrow$ Generate Heritage Site} --- procedurally generates a placeholder heritage site for performance testing.
\end{itemize}

\begin{figure}[h]
\centering
\fbox{\rule{0pt}{7cm}\rule{0.85\textwidth}{0pt}}
\caption{Setup Wizard window inside the Unity Editor. Insert screenshot here.}
\label{fig:wizard}
\end{figure}
